\begin{statement}{A.5.4}{Proved}{The remaining normalization class for a persistent ambient obstruction}
Let $Y$ satisfy the CDP hypotheses, and let $S$ be a connected nonnormal reduced fibre support for which
\[
 H^0(S,\Omega_Y^1|_S\otimes L|_S)\ne0
 \quad\hbox{for every }L\in\Pic(Y).
 \tag{A.27}
\]
There is a normalization component $T_i$ whose minimal resolution $Z_i$ has Kodaira dimension $-\infty$. Every line from $\Pic(Y)$ pulls back to a unitary-flat line on that $Z_i$. A rational or class VII $Z_i$ has trivial Picard pullback from $Y$, and then
\[
 H^2\!\left(T_i,\omega_{T_i}\otimes
             \eta^*((K_Y^{-1}\otimes L)|_S)|_{T_i}\right)\simeq\C
 \quad\hbox{for every }L\in\Pic(Y).
 \tag{A.28}
\]
The same conclusion holds for the exceptional-section case of A.5.3. Thus the only normalization type left by this argument, when (A.28) has not been obtained, is an irrational ruled surface with nonconstant Picard pullback; in base genus at least two its divisor (A.26) has its horizontal section in the strict conductor.
\end{statement}
\begin{proof}
A.4.9 gives $H^0(S,\widetilde\Omega_S^1\otimes L|_S)=0$ on a dense open subset of $\Pic(Y)$. By A.1.1 and (A.27), $H^0(S,K\otimes L|_S)$ is nonzero there. Its nonvanishing locus is closed analytic in the holomorphic Picard family, hence is all of $\Pic(Y)$.

Pull $K$ to each $Z_i$, quotient the resulting sheaf by its torsion, and take its double dual. The result is a fixed holomorphic line bundle $M_i$, since its generic rank on each component is one and $Z_i$ is a smooth surface. The induced maps give an injection
\[
 H^0(S,K\otimes L|_S)\hookrightarrow
       \bigoplus_i H^0(Z_i,M_i\otimes h_i^*L).
 \tag{A.29}
\]
Indeed a section in the kernel vanishes generically on every component, and $K\subset\Omega_Y^1|_S$ is torsion-free over the reduced space. Each summand's nonvanishing locus is closed analytic in the connected complex vector space $\Pic(Y)=H^1(Y,\OO_Y)$. A finite union of proper analytic subsets cannot cover that space. Consequently one component $i$ has
$H^0(Z_i,M_i\otimes h_i^*L)\ne0$ for every $L$.

Fix a Gauduchon metric on $Z_i$, with fundamental form $\omega$ satisfying $\partial\bar\partial\omega=0$. The degree of a line bundle is the integral of its real Chern curvature against $\omega$; it is independent of the chosen Hermitian metric and additive under tensor product. A nonzero section has an effective zero divisor, so its line has nonnegative degree. For any $G\in\Pic(Y)$ the preceding nonvanishing applied to $G^{\otimes n}$ therefore gives
\[
 \deg_\omega(M_i)+n\deg_\omega(h_i^*G)\ge0
 \quad\hbox{for every }n\in\Z.
\]
It follows that $\deg_\omega(h_i^*G)=0$ for every $G$ and every Gauduchon metric.

We record the precise flatness implication. For a topologically trivial line, the real Chern curvature $F$ is exact. If $\int F\wedge\omega=0$, the Gauduchon scalar elliptic equation supplies a real $u$ with $\Lambda_\omega(F-i\partial\bar\partial u)=0$. The remaining form is closed, exact, real and primitive of type $(1,1)$, hence anti-selfdual. Its square integrates both to zero by exactness and to minus its squared norm. It is therefore zero. Changing the line metric by $e^{-u}$ gives zero Chern curvature, hence a flat unitary connection. This is exactly the curvature argument of \cite[Lemma~2.3(1)--(2)]{Teleman2006}. Since $H^2(Y,\Z)=0$, it applies to every $h_i^*G$.

The divisor $D_i$ of A.5.3 is nonzero. For an irreducible nonnormal support this follows from its nonempty conductor; for a connected reducible support every component meets another component, again giving a nonempty conductor. If $Z_i$ had nonnegative Kodaira dimension, some positive pluricanonical power would have a nonzero section, making $\deg_\omega(K_{Z_i})\ge0$. But $D_i$ is a nonzero effective divisor and therefore has strictly positive degree. Equation (A.25), with $\omega_S=(K_Y\otimes\OO_Y(S))|_S$, would give
\[
 0=\deg_\omega h_i^*(K_Y\otimes\OO_Y(S))
   =\deg_\omega(K_{Z_i})+\deg_\omega(D_i)>0,
\]
a contradiction. Thus $\kappa(Z_i)=-\infty$, whose classification consists of rational, irrational ruled, and class VII surfaces.

For a rational $Z_i$, topologically trivial lines are trivial. For a class VII $Z_i$, $h^1(\OO_{Z_i})=1$ and $b_1(Z_i)=1$. With our exponential convention $u\mapsto\exp(2\pi\mathrm i u)$, the unitary-flat subgroup of $\Pic^0(Z_i)$ has real tangent space equal to the injective image of $H^1(Z_i,\mathbb R)$ in $H^1(Z_i,\OO_{Z_i})$. This is \cite[Lemma~2.3(2)]{Teleman2006} after multiplication by the fixed scalar $\mathrm i$, or follows from the pluriharmonic maximum principle. That real tangent space has dimension one. The holomorphic Picard pullback has complex-linear differential and, by the proved flatness, takes values in this real one-dimensional subspace. Its differential is zero, so the homomorphism is constant and equals zero at the trivial line. Every $h_i^*L$ is therefore trivial. Projection formula and $(\pi_i)_*\OO_{Z_i}=\OO_{T_i}$ again descend triviality to $T_i$. The exceptional-section case has the same conclusion by A.5.3.

In each of these cases, the line multiplying $\omega_{T_i}$ in (A.28) is trivial, since it too comes from $Y$. A normal complex surface is Cohen--Macaulay, so its dualizing sheaf and Serre duality give $H^2(T_i,\omega_{T_i})=H^0(T_i,\OO_{T_i})^\vee=\C$. This proves (A.28) and the stated exact reduction of the remaining normalization problem.
\end{proof}


% A.5.5 from part_A6_critical_curves.tex

